\documentclass{article}


% if you need to pass options to natbib, use, e.g.:
%     \PassOptionsToPackage{numbers, compress}{natbib}
% before loading neurips_2024


% ready for submission
\usepackage{neurips_2025}
\usepackage{amsmath}

% to compile a preprint version, e.g., for submission to arXiv, add add the
% [preprint] option:
%     \usepackage[preprint]{neurips_2024}


% to compile a camera-ready version, add the [final] option, e.g.:
%     \usepackage[final]{neurips_2024}


% to avoid loading the natbib package, add option nonatbib:
%    \usepackage[nonatbib]{neurips_2024}


\usepackage[utf8]{inputenc} % allow utf-8 input
\usepackage[T1]{fontenc}    % use 8-bit T1 fonts
\usepackage{hyperref}       % hyperlinks
\usepackage{url}            % simple URL typesetting
\usepackage{booktabs}       % professional-quality tables
\usepackage{amsfonts}       % blackboard math symbols
\usepackage{nicefrac}       % compact symbols for 1/2, etc.
\usepackage{microtype}      % microtypography
\usepackage{xcolor}         % colors



\usepackage{algorithm}
\usepackage{algorithmic}
\usepackage{wrapfig}
\usepackage{graphicx}
\usepackage{enumerate}
\usepackage{multirow}
\usepackage{wrapfig}
\usepackage{lipsum}
\usepackage{xcolor}
\usepackage{diagbox} 
\usepackage{enumitem}
\usepackage{subcaption}


\usepackage{amsmath,amssymb,mathtools}

\newcommand{\R}{\mathbb{R}}
\newcommand{\E}{\mathbb{E}}
\newcommand{\norm}[1]{\left\lVert #1 \right\rVert}
\newcommand{\inner}[2]{\left\langle #1, #2 \right\rangle}
\newcommand{\sigmoid}{\sigma}
\newcommand{\tr}{\mathrm{tr}}
\newcommand{\eps}{\varepsilon}


\title{Manifold-of-Thought: A Geometry-Driven Reasoning Paradigm for Large Language Models}


% The \author macro works with any number of authors. There are two commands
% used to separate the names and addresses of multiple authors: \And and \AND.
%
% Using \And between authors leaves it to LaTeX to determine where to break the
% lines. Using \AND forces a line break at that point. So, if LaTeX puts 3 of 4
% authors names on the first line, and the last on the second line, try using
% \AND instead of \And before the third author name.


\author{%
  David S.~Hippocampus\thanks{Use footnote for providing further information
    about author (webpage, alternative address)---\emph{not} for acknowledging
    funding agencies.} \\
  Department of Computer Science\\
  Cranberry-Lemon University\\
  Pittsburgh, PA 15213 \\
  \texttt{hippo@cs.cranberry-lemon.edu} \\
  % examples of more authors
  % \And
  % Coauthor \\
  % Affiliation \\
  % Address \\
  % \texttt{email} \\
  % \AND
  % Coauthor \\
  % Affiliation \\
  % Address \\
  % \texttt{email} \\
  % \And
  % Coauthor \\
  % Affiliation \\
  % Address \\
  % \texttt{email} \\
  % \And
  % Coauthor \\
  % Affiliation \\
  % Address \\
  % \texttt{email} \\
}


\begin{document}


\maketitle

\begin{abstract}
Large Language Models (LLMs) often benefit from explicit intermediate reasoning, yet existing reasoning methods largely rely on prescribing an external thinking procedure in advance, such as linear chains, parallel sampling, or structured search. While effective in some settings, these approaches treat reasoning as a fixed script rather than an adaptive process, making it difficult to align reasoning behavior with the model's evolving internal logic.

We introduce \textbf{Manifold-of-Thought (MoT)}, a geometry-driven reasoning paradigm that represents reasoning as a continuous logic trajectory in a compact internal state space. Instead of predefining discrete reasoning patterns, MoT maintains a sentence-level state that tracks the evolving logic of the reasoning process as it moves on an internal manifold, and updates this state online as reasoning unfolds. MoT then builds a \emph{training-free logic atlas} from quality-filtered, mixed-dataset reasoning trajectories via two-phase data collection and RBF-weighted aggregation; the resulting head importance profiles are precomputed on a uniform grid in the logic state space, enabling O(1) interpolation for deterministic one-step look-ahead gating during inference, without training any separate policy. To avoid over-intervention in normal reasoning, MoT partitions the logic space into a \emph{safe zone} (minimal or no gating) and a \emph{danger zone} built from pivot steps of incorrect trajectories; suppression is applied only in the danger zone and, when used, conditioned on the current head profile to reduce false positives.

By viewing reasoning through the geometry of its internal logic flow, MoT shifts LLM reasoning from \emph{external script design} to \emph{internal state-space control}. This perspective provides a principled and lightweight \emph{training-free} foundation for adaptive reasoning beyond fixed prompting-based strategies.
\end{abstract}


\section{Introduction}
\label{sec:intro}

\paragraph{Background.}
Reasoning augmentation has become central to improving Large Language Models (LLMs) on complex tasks. Methods such as chain-of-thought prompting, best-of-$N$ sampling, and tree-style search demonstrate that intermediate reasoning can substantially affect final performance. At the same time, they expose a persistent tension: stronger reasoning often comes at the cost of longer generations, greater inference-time computation, and a higher risk of unproductive or looping intermediate steps.

\paragraph{Motivation: reasoning is still externally scripted.}
Despite their differences, many existing reasoning methods share the same structural assumption: the reasoning procedure is chosen \emph{before} inference and then imposed on the model as an external script. Chain-of-thought commits to linear step-by-step derivation; best-of-$N$ commits to parallel exploration followed by selection; tree-based methods commit to branching and backtracking. What is fixed in these approaches is not only the prompt format, but the reasoning procedure itself.

This rigidity matters for two reasons. First, different tasks---and even different stages of the same task---may require different reasoning behaviors, such as exploration, consolidation, verification, or revision. Second, fixed reasoning scripts may continue producing intermediate steps even when the model has already converged to a reliable answer or, conversely, has become stalled and is no longer making meaningful progress. A fixed script therefore cannot adapt reasoning behavior to the model's evolving internal logic.

\paragraph{Key perspective.}
We argue that reasoning should not be viewed as selecting from a small set of discrete thinking scripts. Instead, it should be modeled as the evolution of a compact internal logic state in a continuous geometric space. Under this view, phenomena such as exploration, verification, redirection, and backtracking need not be predefined as explicit reasoning modes; they can instead be understood as different regions, directions, and local motions of an underlying logic trajectory.

The central challenge is therefore to construct a compact state representation that captures the evolving logic of reasoning over time, while remaining sufficiently general to support adaptive control without manually labeling reasoning modes.

\paragraph{Manifold-of-Thought.}
We propose \textbf{Manifold-of-Thought (MoT)}, a geometry-driven reasoning paradigm that shifts LLM reasoning from \emph{external script selection} to \emph{internal state-space control}. MoT operates at sentence boundaries. After each sentence, it summarizes hidden-state dynamics into a compact reasoning state that tracks the current logic of the reasoning process, interprets this process as a point on an evolving internal manifold, and uses the resulting state to control the next reasoning step through head-level soft gating.

This yields a simple closed loop:
\[
\text{observe hidden-state dynamics}
\;\rightarrow\;
\text{update compact reasoning state}
\;\rightarrow\;
\text{control the next reasoning step}.
\]

At the same time, the same evolving logic trajectory determines whether the reasoning process should continue: the loop terminates once the model reaches a sufficiently confident convergent state or, alternatively, exhibits persistent stagnation. In this way, MoT treats reasoning as a continuous control problem whose behavior is governed by the evolving local geometry of the reasoning trajectory, rather than by a predefined prompting procedure.

\paragraph{Why a manifold view?}
Reasoning does not unfold along a single monotonic axis, because its internal logic may advance, hesitate, redirect, revisit earlier lines of thought, or partially reverse course. As a result, a single scalar signal is generally insufficient to characterize reasoning evolution. MoT therefore adopts a compact low-dimensional state whose trajectory captures both the local geometry and the inter-step dynamics of logic flow. We use the term \emph{manifold-of-thought} to emphasize this perspective: reasoning is modeled as motion of internal logic in a compact state space, rather than as jumps among predefined discrete modes or fixed external scripts.

\paragraph{Contributions.}
\begin{itemize}[leftmargin=1.3em, itemsep=0.2em]
    \item \textbf{A geometry-driven formulation of reasoning.} We formulate LLM reasoning as a continuous, controllable trajectory of internal logic in a compact state space, moving beyond externally fixed thinking scripts.
    \item \textbf{A minimal sentence-level instantiation.} We design a simple MoT control loop that summarizes hidden-state dynamics at sentence boundaries and uses a unified compact logic state with deterministic one-step look-ahead retrieval to modulate behavior through head-level soft gating.
    \item \textbf{A training-free adaptive-reasoning framework for frozen LLMs.} MoT builds a mixed-dataset logic atlas through two-phase quality-filtered trajectory collection and RBF-weighted aggregation; head importance profiles are precomputed on a discrete grid in logic state space and queried online via interpolation, requiring no controller training and no gradient updates to the LLM. Safe and danger zones (the latter from pivot steps of incorrect trajectories) restrict intervention to danger regions with optional conditional suppression on the current head profile.
\end{itemize}


\section{Preliminaries}
\label{sec:prelim}

\subsection{Transformer Notation}
Consider an autoregressive transformer LLM with $L$ layers, $H$ attention heads per layer, and hidden size $D$. 
For the $i$-th token at layer $\ell$, denote its hidden state by
\[
h^{(\ell)}_{i} \in \R^{D}.
\]

MoT operates on a \emph{frozen} backbone LLM and adds a lightweight control interface on top of compact sentence-level logic states.

\subsection{Sentence-Level Closed Loop}
We decompose a reasoning trajectory into sentence-level steps
\[
y_{1}, y_{2}, \dots, y_{T},
\]
where each $y_t$ is the text generated between two sentence boundaries (e.g., punctuation, line break, or a predefined delimiter). In practice we use a single consistent definition (e.g., newline or period as delimiter) for both atlas construction and inference. The prompt $x$ is treated as step $0$ only for state initialization.

At the end of step $t$, MoT extracts a compact logic state $m_t$ from the hidden-state dynamics of $y_t$. 
This logic state then controls the generation of the next step $y_{t+1}$. 
Thus, MoT forms a sentence-level closed loop:
\[
y_t \;\longrightarrow\; m_t \;\longrightarrow\; \text{control for } y_{t+1}.
\]

Sentence-level control is a natural compromise between token-level responsiveness and trajectory-level stability: it is fine-grained enough to capture multi-stage reasoning, yet coarse enough to keep the controller stable and interpretable.

\subsection{Adaptive Termination of the Reasoning Loop}
MoT does not fix the number of reasoning steps in advance. Instead, reasoning unfolds sentence by sentence until the evolving logic state indicates either convergence or persistent stagnation. Thus, the reasoning horizon is determined online by the closed loop itself rather than by a predefined script or a fixed-length reasoning budget.


\section{Manifold-of-Thought (MoT)}
\label{sec:method}

\subsection{Overview}
MoT models reasoning as a continuous trajectory of internal logic in a compact state space rather than as a sequence of discrete thinking modes. 
At the end of each sentence-level step $t$, we extract a small set of hidden-state observables that summarize both the current step and its transition from previous steps. 
These observables are normalized (using fixed statistics) into a compact logic state $m_t$, which serves as a lightweight proxy for the evolving logic flow of reasoning.

The same $m_t$ is then used to query a training-free logic atlas and drive a single geometry-driven control mechanism: deterministic one-step look-ahead head-level soft gating for the next reasoning step. The sentence-level loop then terminates adaptively when the logic trajectory indicates convergence or persistent stagnation.

The key design principle is simplicity: MoT does not enumerate discrete reasoning modes such as ``explore'', ``verify'', or ``backtrack'', and does not train a separate reasoning policy network. Instead, such phenomena are treated as different regions or motions of a continuous logic trajectory estimated from trajectory statistics.

% -------------------------
\subsection{Step Observation: Robust Hidden-State Summary}
\label{sec:obs}
At the end of sentence $y_t$, let $n_t$ be its number of tokens. 
For the $i$-th token in $y_t$, we use the last layer's hidden state (or optionally average the last $k$ layers for robustness):
\begin{equation}
h_{t,i} \;=\; h^{(L)}_{t,i},
\qquad h_{t,i}\in\R^D,
\label{eq:lastk}
\end{equation}
i.e., we take the top-layer representation by default ($k=1$); using $h_{t,i} = \frac{1}{k}\sum_{\ell=L-k+1}^{L} h^{(\ell)}_{t,i}$ with $k>1$ is an optional choice.

We summarize the sentence by its center
\begin{equation}
z_t \;=\; \frac{1}{n_t}\sum_{i=1}^{n_t} h_{t,i} \in \R^D,
\label{eq:zt}
\end{equation}
and its within-sentence dispersion
\begin{equation}
\delta_t \;=\; \tr\!\left(
\frac{1}{n_t}\sum_{i=1}^{n_t}(h_{t,i}-z_t)(h_{t,i}-z_t)^\top
\right)
\;=\; \frac{1}{n_t}\sum_{i=1}^{n_t} \norm{h_{t,i}}_2^2 - \norm{z_t}_2^2.
\label{eq:disp}
\end{equation}
Here $\delta_t$ measures how spread out the token representations within the current sentence are; larger values indicate a more internally dispersed or branched step. The second form avoids forming the covariance matrix.

We also use a low-cost uncertainty signal from token logits:
\begin{equation}
H_t \;=\; \frac{1}{n_t}\sum_{i=1}^{n_t} \mathcal{H}(p_{t,i}),
\label{eq:entropy}
\end{equation}
where $p_{t,i}$ is the next-token distribution at token $i$.

For the prompt $x$, we compute $(z_0,\delta_0,H_0)$ in the same way and use them only to initialize the control state.

% -------------------------
\subsection{Compact Logic Signature: Progress, Consistency, and Uncertainty}
\label{sec:logic-signature}
To capture inter-step dynamics, define the sentence displacement
\begin{equation}
\Delta z_t \;=\; z_t - z_{t-1}.
\label{eq:deltaz}
\end{equation}
We then define two simple dynamical quantities.

First, the progress magnitude:
\begin{equation}
v_t \;=\; \norm{\Delta z_t}_2,
\label{eq:speed}
\end{equation}
which measures how much the reasoning state moves from step $t-1$ to step $t$.

Second, the directional consistency:
\begin{equation}
c_t \;=\;
\begin{cases}
1, & t=1,\\[4pt]
\dfrac{\inner{\Delta z_t}{\Delta z_{t-1}}}
{\norm{\Delta z_t}_2 \norm{\Delta z_{t-1}}_2 + \eps}, & t\ge 2,
\end{cases}
\label{eq:consistency}
\end{equation}
where $\eps>0$ is a small constant. 
The value $c_t\in[-1,1]$ indicates whether the trajectory continues in a similar direction ($c_t\approx 1$), turns to a different direction ($c_t\approx 0$), or reverses / backtracks ($c_t<0$).

We form the minimal sentence-level signature
\begin{equation}
s_t \;=\; [\,\delta_t,\; v_t,\; c_t,\; H_t\,]^\top \in \R^4.
\label{eq:st}
\end{equation}
These four quantities are chosen as a minimal sentence-level summary of logic flow: dispersion captures within-step branching, progress measures inter-step advance, consistency reflects directional continuity or redirection, and uncertainty indicates local instability of the current reasoning step.

\paragraph{Fixed normalization.}
To reduce dependence on scale across models and tasks, we normalize each component using fixed statistics. We compute $(\mu,\sigma)\in\R^4$ once on the source dataset (e.g., by collecting signatures $s_t$ over a pass over source prompts) and use the same $(\mu,\sigma)$ for all trajectories in atlas construction and inference:
\begin{equation}
\hat s_t \;=\; \frac{s_t - \mu}{\sigma + \eps},
\label{eq:norm}
\end{equation}
where subtraction and division are elementwise. This keeps the control-state scale stable without per-trajectory state or EMA updates.

% -------------------------
\subsection{Compact Logic State and Manifold View}
\label{sec:mt}
MoT uses the normalized signature as the compact logic state:
\begin{equation}
m_t \;=\; \hat s_t \;\in\; \R^4.
\label{eq:mt}
\end{equation}
Thus, the control state is low-dimensional and continuous, with no learned recurrent memory or temporal smoothing by default; only a one-step finite-difference extrapolation is used for look-ahead retrieval (Sec.~\ref{sec:predictive}). Optionally, temporal smoothing (e.g., $m_t = \beta m_{t-1} + (1-\beta)\hat s_t$ with $\beta \in [0,1)$) can be applied for added stability.

In MoT, the term \emph{manifold-of-thought} refers to the continuous geometry induced by admissible logic states and their trajectories over time, rather than to an explicitly learned nonlinear manifold. Under this view, reasoning is modeled as the evolution of internal logic in a compact geometric state space, not as jumps between predefined discrete modes. 
For visualization only, one may project the trajectory $\{m_t\}$ to $2$D or $3$D post hoc (e.g., by PCA), but such projection is not part of the core method.

% -------------------------
\subsection{Training-Free Logic Atlas from Mixed Datasets}
\label{sec:controller}

MoT keeps the backbone LLM frozen and avoids training any gate predictor. Instead, it builds a \emph{logic atlas} from a mixed dataset
\[
\mathcal{D}_{\mathrm{mix}}=\bigcup_{k=1}^{K}\mathcal{D}_k,
\]
where each $\mathcal{D}_k$ may exhibit different reasoning tendencies, and the union exposes the model to diverse logic flow patterns. Atlas construction proceeds in two phases described below.

\paragraph{Head-importance profile.}
To connect logic states to head control, we compute a teacher-preserving head-importance profile at each trajectory step. The gate interface is defined by
\begin{equation}
\mathrm{AttnOut}^{(\ell)} \;=\; \sum_{h=1}^{H} g^{(\ell,h)} \, O^{(\ell,h)},
\label{eq:headout}
\end{equation}
and head importance is estimated by single-head masking under teacher forcing:
\begin{equation}
r_{t}^{(\ell,h)}
\;=\;
\frac{1}{n_{t+1}}
\sum_{i=1}^{n_{t+1}}
D_{\mathrm{KL}}
\!\left(
p^{\mathrm{full}}_{t+1,i}
\;\|\;
p^{(-\ell,h)}_{t+1,i}
\right),
\label{eq:teacher_importance}
\end{equation}
then normalized within each layer:
\begin{equation}
\bar r_{t}^{(\ell,h)}
\;=\;
\frac{r_{t}^{(\ell,h)}}
{\max_{h'} r_{t}^{(\ell,h')} + \eps}.
\label{eq:teacher_importance_norm}
\end{equation}

\paragraph{Phase 1 -- Bootstrap collection.}
We run teacher rollouts on $\mathcal{D}_{\mathrm{mix}}$ with all gates fixed to $1$ (unperturbed LLM) and collect $N_0$ independent trajectories per problem. From these, we retain only \emph{correct} trajectories and assign each trajectory a quality weight that rewards shorter, more efficient reasoning relative to the best observed length for that problem:
\begin{equation}
q^{(k)}_i \;=\; \exp\!\left(-\alpha\,\frac{L^{(k)}_i - L^*_i}{L^*_i + 1}\right),
\quad L^*_i = \min_{k':\,y^{(k')}_i=1} L^{(k')}_i,
\label{eq:quality}
\end{equation}
where $L^{(k)}_i$ is the number of reasoning sentences in trajectory $k$ of problem $i$, and $L^*_i$ is the length of the shortest correct trajectory. Incorrect trajectories receive zero weight and are discarded. The initial atlas is
\[
\mathcal{A}_0 = \bigl\{\,(m_t,\,d_t,\,\bar r_t,\,q)\;\big|\;(m_t,d_t,\bar r_t)\text{ from correct trajectory step},\;q\text{ by Eq.~}\eqref{eq:quality}\bigr\}.
\]

\paragraph{Phase 2 -- On-policy refinement.}
Because bootstrap trajectories are collected under gate $=1$, the resulting logic-state distribution may differ from the distribution encountered during gated inference. We therefore iteratively refine the atlas using \emph{on-policy} rollouts. At iteration $n\ge 1$, we collect one trajectory per problem using the gate predictor $G_n$ built from the current atlas $\mathcal{A}_{n-1}$, filter for correctness and quality as in Phase~1, and merge the new points into the atlas:
\[
\mathcal{A}_n \;=\; \mathcal{A}_{n-1} \;\cup\; \mathcal{A}^{\mathrm{new}}_n.
\]
We monitor convergence via two complementary criteria measured between consecutive iterations. Let $P_n(m)$ denote the empirical logic-state distribution of $\mathcal{A}_n$. The \emph{distribution shift} is estimated by a binned KL divergence on the 4D logic space:
\begin{equation}
D^{(n)}_{\mathrm{KL}} \;=\; D_{\mathrm{KL}}\!\bigl(P_n(m)\,\|\,P_{n-1}(m)\bigr),
\label{eq:kl_conv}
\end{equation}
and the \emph{gate-profile shift} measures how much the predicted head-importance changes across the atlas:
\begin{equation}
\Delta^{(n)}_{\mathrm{gate}} \;=\; \E_{m\sim\mathcal{A}_n}\!\left[\bigl\|G_n(m) - G_{n-1}(m)\bigr\|_2\right].
\label{eq:gate_conv}
\end{equation}
Refinement stops when both quantities fall below fixed thresholds $\epsilon_{\mathrm{KL}}$ and $\epsilon_{\mathrm{gate}}$, or after a maximum number of iterations. In practice, convergence is reached in 2--3 rounds.

\paragraph{RBF-weighted atlas and offline grid precomputation.}
Once the final atlas $\mathcal{A}^*$ is available, we define the gate-prediction map by quality-weighted RBF aggregation over the directional partition $d\in\{+,-\}$:
\begin{equation}
G(m,d) \;\triangleq\;
\frac{\displaystyle\sum_{j:\,d_j=d}\; q_j\,\kappa(m,m_j)\,\bar r_j}
     {\displaystyle\sum_{j:\,d_j=d}\; q_j\,\kappa(m,m_j)+\eps},
\quad
\kappa(m,m_j)=\exp\!\left(-\frac{\|m-m_j\|^2}{2\sigma^2}\right),
\label{eq:rbf}
\end{equation}
where $d_j$ is the directional partition index of atlas point $j$ (from the sign of $c_t$ at that step), and $\sigma$ is set to the median pairwise distance within each partition. The directional partition index $d_t\in\{+,-\}$ is defined by
\[
d_t =
\begin{cases}
+, & c_t \ge 0,\\
-, & c_t < 0,
\end{cases}
\]
and acts as a geometric retrieval index that separates forward-progress neighborhoods from redirection or backtracking neighborhoods, reducing multimodal interference in the aggregation. It is not a symbolic logic-mode label.

Because the logic state $m_t\in\R^4$ is low-dimensional, we discretize the logic space into a uniform grid (e.g.\ $B$ bins per dimension, yielding $B^4$ cells) and \emph{precompute} $G(m,d)$ at every grid center offline. During online inference, querying the map reduces to a constant-time 4D multilinear interpolation over the 16 surrounding grid vertices:
\begin{equation}
G(m,d) \;\approx\; \sum_{\mathbf{u}\,\in\,\mathrm{vertices}(m)} \alpha_{\mathbf{u}}(m)\cdot G(m_{\mathbf{u}},d),
\label{eq:interp}
\end{equation}
where $\alpha_{\mathbf{u}}(m)\ge 0$ are the standard 4D multilinear interpolation weights (summing to $1$), and $m_{\mathbf{u}}$ are the $2^4=16$ grid vertex positions surrounding $m$. This makes online gate retrieval an O(1) memory lookup with negligible latency overhead, while preserving the smooth, continuous variation of gate weights as the logic state evolves through the space.

% -------------------------
\paragraph{Safe zone and danger zone (negative signal).}
The atlas above is built only from \emph{correct} trajectories, so it describes ``where good reasoning tends to use heads'' but not ``where reasoning goes wrong.'' To avoid over-intervening in regions where the model is already reasoning normally (and to avoid the accuracy drop observed when applying a global correct-atlas blend), we partition the logic state space into a \textbf{safe zone} and a \textbf{danger zone}. The safe zone comprises states visited by correct trajectories and the large overlap region where both correct and incorrect trajectories pass; the danger zone comprises states where \emph{incorrect} trajectories exhibit a \emph{divergence} in head usage relative to the correct atlas---the ``pivot'' steps where the trajectory goes off track.

We define a \textbf{pivot} as follows. For each \emph{incorrect} trajectory, we compute head-importance profiles $r_t$ at each step (same as for correct trajectories). At step $t$, let $\bar r_{\mathrm{correct}}(m_t,d_t) = G(m_t,d_t)$ be the correct-atlas query. The \emph{deviation} is
\begin{equation}
\mathrm{dev}_t \;=\; \bigl\| r_t - \bar r_{\mathrm{correct}}(m_t,d_t) \bigr\|_2.
\label{eq:deviation}
\end{equation}
The \textbf{pivot} of that trajectory is the step with the largest $\mathrm{dev}_t$ (or the first step above a threshold $\tau_{\mathrm{dev}}$). Each pivot yields a state $(m_{\mathrm{pivot}}, d_{\mathrm{pivot}})$, a head profile $r_{\mathrm{pivot}}$, and the correct-atlas value $\bar r_{\mathrm{correct}}$ at that state. The \textbf{divergence vector} is
\begin{equation}
\Delta r \;=\; r_{\mathrm{pivot}} - \bar r_{\mathrm{correct}}.
\label{eq:delta_r}
\end{equation}
We build a \textbf{danger map} $\Delta r(m,d)$ by RBF interpolation over pivot points (or a precomputed grid), so that near pivot states we know ``which head pattern to suppress.''

\paragraph{Intervention policy: minimal in safe zone, suppression in danger zone.}
In the \textbf{safe zone}, we apply \emph{minimal or no intervention}: blended gates are kept close to the ungated model ($\lambda_{\mathrm{safe}}\approx 0$ or gates $=1$), so the model reasons naturally. In the \textbf{danger zone}, we apply \emph{suppression} via Eq.~\eqref{eq:g_apply_danger}. The default implementation determines safe vs.\ danger by whether the RBF-interpolated \emph{strength} at $(m_t,d_t)$ (i.e., proximity to pivot points) exceeds a threshold: if strength is below the threshold, the state is treated as safe; otherwise it is treated as danger and suppression is applied. Optionally, one can add \emph{conditional suppression}: compute the current step's $r_t$ and $\mathrm{deviation}_{\mathrm{now}} = \|r_t - \bar r_{\mathrm{correct}}(m_t,d_t)\|_2$, and apply suppression only if $\mathrm{deviation}_{\mathrm{now}} > \tau_{\mathrm{dev}}$. When suppression is applied (by strength or by the optional deviation check), we set
\begin{equation}
g_{\mathrm{apply}}^{(\ell,h)} \;=\; \mathrm{clip}\bigl( \bar r_{\mathrm{correct}}^{(\ell,h)} - \alpha \cdot \Delta r^{(\ell,h)}(m_t,d_t),\;\; 0,\; 1 \bigr),
\label{eq:g_apply_danger}
\end{equation}
and blend with the ungated model using $\lambda_{\mathrm{danger}}$. This way, the same $(m,d)$ can be ``safe'' when the model's current $r_t$ matches the correct pattern, and ``danger'' when $r_t$ deviates, avoiding false-positive intervention on correct trajectories that pass through the same region.

% -------------------------
\subsection{Deterministic Look-Ahead Atlas Gating}
\label{sec:predictive}
At sentence step $t$, MoT uses a deterministic one-step look-ahead query to avoid lagging behind local logic transitions:
\begin{equation}
\hat m^{\mathrm{lite}}_{t+1} \;=\; m_t + \rho_t\,(m_t-m_{t-1}),
\label{eq:predict_state_lite}
\end{equation}
with
\begin{equation}
\rho_t \;=\; \frac{\rho_0}{1+\alpha_H H_t + \alpha_c(1-c_t)},
\qquad \rho_t\in[\rho_{\min},\rho_{\max}].
\label{eq:rho}
\end{equation}
Let $d_t$ be the directional partition index determined by the sign of $c_t$. We then retrieve the correct-atlas profile
\begin{equation}
\tilde g^{(\ell,h)}_{t+1}
\;=\;
G^{(\ell,h)}(\hat m^{\mathrm{lite}}_{t+1}, d_t)\in[0,1].
\label{eq:gh_main}
\end{equation}

\paragraph{Safe vs danger gating.}
If the current state $(\hat m^{\mathrm{lite}}_{t+1}, d_t)$ is in the \textbf{safe zone} (or danger map returns negligible $\Delta r$), we apply minimal blending:
\begin{equation}
g^{(\ell,h)}_{t+1} \;=\; (1-\lambda_{\mathrm{safe}})\cdot 1 + \lambda_{\mathrm{safe}}\, \tilde g^{(\ell,h)}_{t+1},
\quad \lambda_{\mathrm{safe}}\approx 0.
\label{eq:blend_safe}
\end{equation}
If the state is in the \textbf{danger zone} (strength above threshold, or when using optional conditional suppression and $\mathrm{deviation}_{\mathrm{now}} > \tau_{\mathrm{dev}}$), we use $g_{\mathrm{apply}}$ from Eq.~\eqref{eq:g_apply_danger} instead of $\tilde g$:
\begin{equation}
g^{(\ell,h)}_{t+1} \;=\; (1-\lambda_{\mathrm{danger}})\cdot 1 + \lambda_{\mathrm{danger}}\, g_{\mathrm{apply}}^{(\ell,h)}.
\label{eq:blend_danger}
\end{equation}
This keeps the controller single-path and implementation-stable while restricting intervention to danger regions and only when the head profile looks wrong.

% -------------------------
\subsection{Adaptive Termination Within the Closed Loop}
\label{sec:termination}
MoT does not treat stopping as a separate control module. Instead, termination is part of the sentence-level reasoning loop itself. We apply \emph{any} stopping check only after at least $t_{\min}$ reasoning steps (default $t_{\min}=2$), so that both conditions below are active only once the trajectory has progressed beyond the initial steps.

\paragraph{Convergence.}
We regard the reasoning process as converged when uncertainty is sufficiently low, progress becomes small, and the trajectory remains directionally stable:
\begin{equation}
H_t < \tau_H,\qquad v_t < \tau_v,\qquad c_t > \tau_c.
\label{eq:converge}
\end{equation}
This corresponds to a state in which the model appears to have reached a stable logical conclusion.

\paragraph{Persistent stagnation.}
We also stop when the reasoning process ceases to make meaningful progress:
\begin{equation}
v_t < \tau_{\mathrm{prog}}
\quad \text{for } r_{\mathrm{stop}} \text{ consecutive steps}.
\label{eq:stagnation}
\end{equation}
To avoid early stop on benign slowdowns, a practical variant additionally requires non-improving uncertainty over the same window. The stagnation criterion is likewise applied only when $t \geq t_{\min}$.

Together, these rules allow MoT to terminate reasoning adaptively without introducing an additional resource-control module beyond soft gating.

% -------------------------
\subsection{MoT Inference Algorithm}
\label{sec:algo}

Algorithm~\ref{alg:mot} summarises the full online inference loop. The only offline cost is atlas construction (Sec.~\ref{sec:controller}); the precomputed grid requires no gradient updates and is built once for a given backbone.

\begin{algorithm}[t]
\caption{MoT Inference (Sentence-Level Closed Loop)}
\label{alg:mot}
\begin{algorithmic}[1]
\STATE \textbf{Input:} prompt $x$, frozen LLM with $L$ layers and $H$ heads
\STATE \textbf{Precomputed:} logic atlas grid $\{G(m_{\mathbf{u}}, d)\}$ for all grid vertices $m_{\mathbf{u}}$ and $d\in\{+,-\}$,\\
\hspace{4.2em} fixed normalization statistics $(\mu,\sigma)$
\STATE Treat $x$ as step $0$; compute $(z_0,\delta_0,H_0)$ from prompt hidden states
\STATE Set $v_0=0$, $c_0=1$, form $s_0$, normalize to $\hat s_0$ via Eq.~\eqref{eq:norm}, set $m_0 = \hat s_0$
\STATE Set $m_{-1}=m_0$ (initialization for first-step look-ahead)
\FOR{$t=0,1,\dots,T_{\max}-1$}
    \STATE Compute $\rho_t$ via Eq.~\eqref{eq:rho}; form $\hat m^{\mathrm{lite}}_{t+1}$ via Eq.~\eqref{eq:predict_state_lite}
    \STATE Determine directional partition index $d_t = \mathrm{sign}(c_t)$
    \STATE Retrieve $\{\tilde g^{(\ell,h)}_{t+1}\}$ from correct atlas (Eq.~\eqref{eq:interp}, \eqref{eq:gh_main})
    \STATE If in safe zone: blend via Eq.~\eqref{eq:blend_safe} ($\lambda_{\mathrm{safe}}\approx 0$). If in danger zone and (optional) $\mathrm{deviation}_{\mathrm{now}}>\tau_{\mathrm{dev}}$: compute $g_{\mathrm{apply}}$ via Eq.~\eqref{eq:g_apply_danger} and blend via Eq.~\eqref{eq:blend_danger}
    \STATE Generate sentence $y_{t+1}$ from prefix $(x, y_{1:t})$ with gated attention heads
    \STATE Extract hidden states of $y_{t+1}$ via Eq.~\eqref{eq:lastk}; compute $(z_{t+1},\delta_{t+1},H_{t+1})$
    \STATE Compute $\Delta z_{t+1}$; obtain $(v_{t+1},c_{t+1})$ via Eq.~\eqref{eq:speed}--\eqref{eq:consistency}
    \STATE Form $s_{t+1}$ via Eq.~\eqref{eq:st}; normalize to $\hat s_{t+1}$ via Eq.~\eqref{eq:norm}; set $m_{t+1}=\hat s_{t+1}$
    \IF{$t+1\ge t_{\min}$ \textbf{and} (convergence Eq.~\eqref{eq:converge} \textbf{or} stagnation Eq.~\eqref{eq:stagnation})}
        \STATE Decode final answer from current prefix $(x, y_{1:t+1})$; \textbf{break}
    \ENDIF
\ENDFOR
\STATE \textbf{if} no stopping condition met by $T_{\max}$: decode final answer from current prefix
\end{algorithmic}
\end{algorithm}

\section{Logic Flow Without Discrete Modes}
\label{sec:humanlike}

MoT does not classify reasoning steps into discrete templates (e.g., ``brainstorm'', ``verify'', or ``backtrack''). Instead, such phenomena emerge as characteristic regions and motions of the compact logic trajectory induced by $(\delta_t, v_t, c_t, H_t)$ and the logic state $m_t$:
\begin{itemize}[leftmargin=1.3em, itemsep=0.2em]
    \item \textbf{Exploration / information gathering:} elevated dispersion $\delta_t$ and uncertainty $H_t$, often indicating that the current step is internally branched and not yet consolidated.
    \item \textbf{Exploitation / convergence:} decreasing uncertainty $H_t$ together with reduced dispersion and strong directional consistency ($c_t \approx 1$), indicating more stable logical progress.
    \item \textbf{Backtracking / strategy shift:} drops in directional consistency, especially $c_t \le 0$, often with non-monotonic uncertainty, indicating redirection or partial reversal in representation space.
    \item \textbf{Stagnation / looping:} persistently small progress $v_t$, often accompanied by repetitive generation and non-improving uncertainty, indicating limited logical advance.
\end{itemize}
These patterns are not imposed as symbolic mode labels; they arise naturally from the geometry of the evolving logic state. We examine these correspondences through trajectory visualizations, stopping-state analyses, and gate analyses in the interpretability study.
 
\section{Experiments}
\label{sec:experiments}

\subsection{Evaluation Goals}
Experiments validate three claims aligned with our motivation:
\begin{enumerate}[leftmargin=1.3em, itemsep=0.2em]
    \item \textbf{Paradigm / adaptivity:} MoT adapts reasoning behavior across stages and tasks without discrete mode definitions.
    \item \textbf{Efficiency:} MoT reduces unnecessary reasoning by stopping once the logic trajectory indicates convergence or persistent stagnation.
    \item \textbf{Quality / robustness:} MoT maintains or improves accuracy through state-aware soft gating while avoiding unproductive loops.
\end{enumerate}

\subsection{Tasks and Benchmarks}
We evaluate MoT on three benchmark families that stress complementary aspects of reasoning:
\begin{itemize}[leftmargin=1.3em, itemsep=0.2em]
    \item \textbf{Sequential reasoning:} problems requiring multi-step derivation, where adaptive logical progression is critical.
    \item \textbf{Revision-sensitive reasoning:} problems where incorrect intermediate assumptions must be revised, making logic redirection and backtracking especially important.
    \item \textbf{Long-horizon reasoning:} problems prone to overthinking or looping, where adaptive stopping is especially important.
\end{itemize}
We instantiate each family with standard benchmarks and keep model, prompting, and evaluation protocols consistent across methods.

\subsection{Experimental Setup}
The backbone LLM is frozen in all experiments. We build a training-free logic atlas from mixed source datasets following the two-phase protocol in Sec.~\ref{sec:controller}: Phase~1 collects quality-filtered trajectories under unperturbed inference (gate$=1$); Phase~2 performs on-policy refinement until the KL divergence $D^{(n)}_{\mathrm{KL}}$ and gate-profile shift $\Delta^{(n)}_{\mathrm{gate}}$ both fall below their convergence thresholds (Eq.~\eqref{eq:kl_conv}--\eqref{eq:gate_conv}), typically within 2--3 rounds. The final atlas is discretized onto a uniform $B^4$ grid using RBF-weighted aggregation (Eq.~\eqref{eq:rbf}), and all online gate queries are answered by 4D multilinear interpolation (Eq.~\eqref{eq:interp}) with O(1) cost. No gradient-based optimization is used at any stage. Unless otherwise stated, a single atlas is built once and transferred unchanged to all target evaluation datasets. A recommended transfer protocol is: build atlas on source datasets, then evaluate on the same plus additional unseen target datasets without any atlas updates.

Stopping thresholds are selected on a held-out split of the source mixture without requiring ground-truth intermediate traces, as follows. We enumerate candidate threshold configurations and retain those under which average reasoning length is minimized while final accuracy does not fall significantly below the standard CoT baseline on the held-out split. The selected thresholds are then fixed for all target evaluation benchmarks. This procedure prioritizes efficiency while ensuring the trajectory-based stopping does not degrade the model's reasoning quality.

\subsection{Baselines}
\begin{itemize}[leftmargin=1.3em, itemsep=0.2em]
    \item \textbf{Vanilla decoding:} direct answer without explicit intermediate reasoning.
    \item \textbf{Standard CoT:} sentence-by-sentence reasoning with no geometry-driven control.
    \item \textbf{CoT + adaptive stopping only:} standard reasoning with the same stopping rule but without soft gating, isolating the effect of adaptive termination.
    \item \textbf{CoT + atlas gating only:} geometry-driven soft gating from atlas queries, together with a fixed reasoning horizon or default stopping policy, isolating the effect of behavioral modulation.
    \item \textbf{Length-matched CoT:} standard CoT truncated to match the average reasoning length of MoT, testing whether gains come from better control rather than simply shorter reasoning traces.
    \item \textbf{MoT full:} atlas-based head-level soft gating with deterministic look-ahead retrieval, together with adaptive termination inside the same compact logic-state loop.
\end{itemize}

\subsection{Metrics and Pareto Frontiers}
We report both quality and cost, emphasizing Pareto trade-offs:
\begin{itemize}[leftmargin=1.3em, itemsep=0.2em]
    \item \textbf{Quality:} accuracy, exact match, or task-specific success rate.
    \item \textbf{Cost:} (i) total generated tokens, (ii) number of reasoning sentences, and optionally (iii) wall-clock latency.
\end{itemize}
In addition to Pareto frontiers, we compare methods under matched inference budgets to test whether adaptive control yields better quality-cost trade-offs than fixed strategies.

\subsection{Ablations}
Key ablations include:
\begin{itemize}[leftmargin=1.3em, itemsep=0.2em]
    \item \textbf{Vector $m_t$ vs.\ scalar state:} test whether multi-dimensional logic states better capture stage structure than a single control variable.
    \item \textbf{State dimension:} compare minimal 3D/4D variants (e.g., $(\delta,v,H)$ vs.\ $(\delta,v,c,H)$) to test trade-offs between simplicity and control precision.
    \item \textbf{Temporal smoothing:} compare default (no smoothing, $m_t = \hat s_t$) with optional smoothing ($m_t = \beta m_{t-1} + (1-\beta)\hat s_t$, $\beta>0$) to test the value of trajectory-level state tracking.
    \item \textbf{Remove head gates:} stopping only vs.\ MoT full to quantify geometry-driven behavioral control.
    \item \textbf{Remove adaptive stopping:} gating only vs.\ MoT full to quantify the contribution of adaptive horizon control.
    \item \textbf{Atlas ablations:} (i) current-state retrieval $G(m_t,d_t)$ vs.\ deterministic look-ahead retrieval $G(\hat m^{\mathrm{lite}}_{t+1},d_t)$ via Eq.~\eqref{eq:interp}; (ii) partition-unaware aggregation vs.\ directional-partitioned RBF aggregation; (iii) fixed vs.\ adaptive $\rho_t$ schedule; (iv) single-phase bootstrap-only atlas vs.\ two-phase atlas with on-policy refinement.
    \item \textbf{Feature ablations:} remove one component from $(\delta_t, v_t, c_t, H_t)$ at a time to test which aspects of logic flow matter most.
\end{itemize}
Sensitivity to hidden-state summarization (e.g., last-1 vs.\ last-4 average) can be reported in the appendix.

\subsection{Interpretability: Logic Trajectories, Stopping States, and Gate Analyses}
We complement the main experiments with interpretability analyses:
\begin{enumerate}[leftmargin=1.3em, itemsep=0.2em]
    \item \textbf{Logic trajectory visualizations:} project $\hat s_t$ or $m_t$ to 2D and visualize reasoning trajectories across tasks and methods.
    \item \textbf{Event overlays:} mark regions with high $\delta_t$, high $H_t$, low $c_t$, or small $v_t$ to interpret exploration, uncertainty, redirection, and stagnation.
    \item \textbf{Stopping-state analyses:} visualize where trajectories terminate, distinguishing convergent stopping from persistent stagnation.
    \item \textbf{Gate analyses:} visualize how head gates vary over regions of the compact logic state, showing continuous behavioral modulation rather than discrete switching.
\end{enumerate}


\section{Discussion and Potential Limitations}
\label{sec:discussion}

\paragraph{Simplicity and parameter dependence.}
MoT uses a minimal set of observables and a single gate interface, with scale sensitivity reduced by fixed normalization. The atlas is built entirely from trajectory statistics via quality-weighted RBF aggregation and offline grid precomputation; no auxiliary neural controller is trained, and no gradient updates touch the backbone LLM. This keeps the method simple and focused: the present paper isolates the geometry-driven reasoning paradigm itself rather than combining it with multiple auxiliary resource-control modules.

\paragraph{When does MoT help most?}
MoT is most beneficial when reasoning is long, stage-dependent, or prone to looping. In such cases, geometry-driven gating can adapt behavior across reasoning stages, while adaptive termination can prevent unnecessary continuation once the model has either converged or ceased to make meaningful progress.

\paragraph{Limitations.}
The current formulation focuses on sentence-level logic tracking and head-level behavioral modulation. It does not yet address finer-grained token-level control or richer context-management mechanisms. Atlas quality depends on trajectory coverage across the source mixture and on how faithfully the four geometric features $(\delta_t, v_t, c_t, H_t)$ reflect the evolving reasoning process; sparse regions of the logic space may yield less reliable gate estimates, mitigated in practice by the RBF interpolation fallback to neighboring atlas points.

\section{Future Work}
\label{sec:future}

A natural next step is to extend the compact logic-state framework beyond behavioral control. In particular, the same geometric view may support adaptive context selection or memory-aware reasoning, where historical intermediate steps are retained or discarded according to their relation to the evolving logic trajectory. We leave such extensions for future work in order to keep the present paper focused on the core paradigm: geometry-driven reasoning through a single unified soft-gating mechanism.



\bibliographystyle{plain} %%%%% 仅显示数字
% \bibliographystyle{unsrtnat} %%%%% 显示作者信息
\bibliography{reference}




\end{document}